\subsection{nml\_testcase\_tcf}
\label{SEC:nml_testcase_tcf}
\inp*[nmlTestcaseTCF]{}[nml\_testcase\_tcf][]

This \hyperref[SEC:User_Input]{namelist block}
is used to specify inputs for the Taylor-Couette flow test case.
%
This version of the Taylor-Couette flow case is based on that found
in \linecite{Wang_2021}, which uses the initial conditions from the Taylor-Green
vortex problem as inspiration for the initial conditions for this case to speed
up the time it takes for the flow to become fully turbulent.
%
The default values in this namelist determining the domain extents are based on
that paper.
%
The default cell counts here are slightly coarser than their work which used
64, 20, and 80 grid cells in the respective $x$, $r$, and $\theta$ directions.

\begin{nmlblocklst}[
,label={LST:nml_testcase_tcf}
,caption={Namelist nml\_testcase\_tcf.}
]
&nml_testcase_tcf
  Inner_Radius       = 0.5     ! [m]
  Inner_Wall_Spacing = 0.002   ! [m]
  Outer_Radius       = 1.0     ! [m]
  Outer_Wall_Spacing = 0.0025  ! [m]
  Aspect_Ratio       = 6.2832  ! 2*pi
  Axial_Cells        = 24
  Radial_Cells       = 20
  Azimuthal_Cells    = 48
  Theta_Shift        = 0.0     ! [deg]
  Periodic_Sectors   = 1
/
\end{nmlblocklst}

%==========================================================================

\begin{nmldescitem}
Inner_Radius = 0.5  ! [m]
\end{nmldescitem}
\inp*[nmlTestcaseTCF::InnerRadius]{}[Inner\_Radius][]

\begin{nmldescbody}
This gives the radius of the inner wall for the Taylor-Couette flow problem.
\end{nmldescbody}

%==========================================================================

\begin{nmldescitem}
Inner_Wall_Spacing = 0.002  ! [m]
\end{nmldescitem}
\inp*[nmlTestcaseTCF::InnerWallSpacing]{}[Inner\_Wall\_Spacing][]

\begin{nmldescbody}
This gives the wall spacing at the inner wall for the Taylor-Couette flow
problem.
\end{nmldescbody}

%==========================================================================

\begin{nmldescitem}
Outer_Radius = 1.0  ! [m]
\end{nmldescitem}
\inp*[nmlTestcaseTCF::OuterRadius]{}[Outer\_Radius][]

\begin{nmldescbody}
This gives the radius of the outer wall for the Taylor-Couette flow problem.
\end{nmldescbody}

%==========================================================================

\begin{nmldescitem}
Outer_Wall_Spacing = 0.0025  ! [m]
\end{nmldescitem}
\inp*[nmlTestcaseTCF::OuterWallSpacing]{}[Outer\_Wall\_Spacing][]

\begin{nmldescbody}
This gives the wall spacing at the outer wall for the Taylor-Couette flow
problem.
\end{nmldescbody}

%==========================================================================

\begin{nmldescitem}
Aspect_Ratio = 6.2832  ! 2*pi
\end{nmldescitem}
\inp*[nmlTestcaseTCF::AspectRatio]{}[Aspect\_Ratio][]

\begin{nmldescbody}
This is used to specify the axial length, $l_x$, of the domain for the
Taylor-Couette flow problem. The length is determined by the equation:

$
  l_x = \inptext{Aspect\_Ratio} \times
        \left| \inptext{Outer\_Radius}
             - \inptext{Inner\_Radius}
        \right|
$
\end{nmldescbody}

%==========================================================================

\begin{nmldescitem}
Axial_Cells = 24
\end{nmldescitem}
\inp*[nmlTestcaseTCF::AxialCells]{}[Axial\_Cells][]

\begin{nmldescbody}
This specifies the number of axial cells to use for the Taylor-Couette flow
problem.
\end{nmldescbody}

%==========================================================================

\FINALFIX{\newpage}
\begin{nmldescitem}
Radial_Cells = 20
\end{nmldescitem}
\inp*[nmlTestcaseTCF::RadialCells]{}[Radial\_Cells][]

\begin{nmldescbody}
This specifies the number of radial cells to use for the Taylor-Couette flow
problem.
\end{nmldescbody}

%==========================================================================

\begin{nmldescitem}
Azimuthal_Cells = 48
\end{nmldescitem}
\inp*[nmlTestcaseTCF::AzimuthalCells]{}[Azimuthal\_Cells][]

\begin{nmldescbody}
This specifies the number of azimuthal cells to use for the Taylor-Couette
flow problem.
\end{nmldescbody}

%==========================================================================

\begin{nmldescitem}
Theta_Shift = 0.0  ! [deg]
\end{nmldescitem}
\inp*[nmlTestcaseTCF::ThetaShift]{}[Theta\_Shift][]

\begin{nmldescbody}
This specifies an angle (in degrees) to shift the grid in theta for the
Taylor-Couette flow problem. The purpose of this input is to rotate the grid
so that a cell boundary is not located at the angle 0\slash 360 degrees.
\end{nmldescbody}

%==========================================================================

\begin{nmldescitem}
Periodic_Sectors = 1
\end{nmldescitem}
\inp*[nmlTestcaseTCF::PeriodicSectors]{}[Periodic\_Sectors][]

\begin{nmldescbody}
This input can be used to run only a sector of the Taylor-Couette flow
problem. This value specifies how many pie slices to divide the grid into,
with the final grid domain consisting of only one pie slice.
\end{nmldescbody}

%==========================================================================

